\section*{Capítulo \ref{ch:b3:rna-regulation} --- ARN no codificantes y regulación postranscripcional}

\begin{solution}{exo:b3:rna-regulation:1}
miARN: codificado en el genoma como una horquilla, \qty{22}{nt},
cortado por \omterm{def:b3:rna-regulation:mirna}{Drosha} y después por \omterm{def:b3:rna-regulation:rnai}{Dicer}, se aparea de forma imperfecta
por su semilla y reprime la traducción y la estabilidad de muchos
mensajeros. \omterm{def:b3:rna-regulation:ncrna}{ARNip}: procede de un ARN largo de doble hebra (vírico, de un
transgén, experimental), \qty{21}{nt}, depende de \omterm{def:b3:rna-regulation:rnai}{Dicer}, se aparea
perfectamente y dirige a \omterm{def:b3:rna-regulation:rnai}{Argonauta} para que corte la diana. \omterm{def:b3:rna-regulation:ncrna}{ARNpi}:
procede de transcritos de hebra sencilla de los \omterm{def:b3:rna-regulation:pirna}{agrupamientos de ARNpi},
\qtyrange{26}{31}{nt}, independiente de \omterm{def:b3:rna-regulation:rnai}{Dicer}, cargado en proteínas PIWI
de la línea germinal, corta los transcritos de transposones (ping-pong) y
dirige el silenciamiento de la \omterm{def:b3:chromatin-epigenetics:nucleosome}{cromatina} del ADN del transposón.
\end{solution}

\begin{solution}{exo:b3:rna-regulation:2}
La Pol II transcribe el pri-miARN; \omterm{def:b3:rna-regulation:mirna}{Drosha} (con DGCR8) recorta la
horquilla en el núcleo; la exportina-5 exporta el pre-miARN; \omterm{def:b3:rna-regulation:rnai}{Dicer} corta
el bucle y deja un dúplex de \qty{22}{nt}; una hebra se carga en
\omterm{def:b3:rna-regulation:rnai}{Argonauta}; la semilla (nucleótidos 2--8) se aparea con un sitio de la
3$'$ UTR; \omterm{def:b3:rna-regulation:rnai}{Argonauta} recluta la maquinaria de desadenilación y
descaperuzado y bloquea la iniciación --- el mensajero se traduce menos y
decae más deprisa.
\end{solution}

\begin{solution}{exo:b3:rna-regulation:3}
El ARN fabricado in vitro por una polimerasa de fago contenía
contaminantes de doble hebra (la polimerasa copia también la otra hebra
del molde y se sale por los extremos), de modo que las preparaciones ``con
sentido'' llevaban el verdadero desencadenante. Fire y Mello purificaron
cada hebra, mostraron que cualquiera de ellas sola hacía poco y luego las
hibridaron a propósito: el ARN de doble hebra era al menos cien veces más
potente, y bastaban unas pocas moléculas por célula.
\end{solution}

\begin{solution}{exo:b3:rna-regulation:4}
Con hambre de hierro: la IRP une los IRE. Ferritina: se bloquea la
iniciación de la traducción y no se fabrica proteína de almacenamiento
(control de la traducción). Receptor de transferrina: el mensajero queda
protegido de su nucleasa y se acumula, de modo que se fabrica más
receptor y se importa más hierro (control de la estabilidad del
mensajero).
\end{solution}

\begin{solution}{exo:b3:rna-regulation:5}
$\gamma = \ln 2/\qty{120}{min} = \qty{0.0058}{min^{-1}}$; $m^{*} =
5/0.0058 \approx 870$ moléculas. Con la degradación duplicada:
$m^{*} \approx 430$, semitiempo $\qty{120}{min}/2 = \qty{60}{min}$.
\end{solution}

\begin{solution}{exo:b3:rna-regulation:6}
Secuencia 3$'$ UTR total, $2\times 10^{7}$ nt; coincidencias al azar $2\times
10^{7}/4^{7} = 2\times 10^{7}/\num{16384} \approx \num{1200}$ sitios. Una diana real se distingue por la conservación del sitio entre
especies, por un contexto favorable (una vecindad rica en AU, un sitio
alejado del codón de parada, apareamiento adicional en el extremo 3$'$
del miARN), por la coexpresión del miARN y de la diana en la misma célula
y por el experimento: el mensajero baja cuando el miARN está presente y
deja de bajar cuando se muta el sitio.
\end{solution}

\begin{solution}{exo:b3:rna-regulation:7}
(a) XX, \emph{tra} nulo: se fabrica SXL pero no TRA, y \emph{dsx} toma el
corte masculino por defecto --- un macho somático (estéril). (b) XY con
TRA constitutiva: TRA-2 está presente en ambos sexos, de modo que DSX se
fabrica en su forma femenina --- desarrollo somático femenino (una
``seudohembra'' estéril, con la línea germinal siguiendo otras señales).
(c) Sin \emph{dsx}: no hay ninguna de las dos formas de DSX, de modo que
no se impone ninguno de los dos programas de diferenciación terminal ---
un intersexo con caracteres mezclados o incompletos.
\end{solution}

\begin{solution}{exo:b3:rna-regulation:8}
En el estado estacionario el mensajero, y por tanto la proteína, es
proporcional a la semivida del mensajero:
$\qty{180}{min}/\qty{15}{min} = 12$ veces más proteína. La proteína es
una señal de crecimiento normal que debe ser transitoria, un pulso tras
cada estímulo; fabricada doce veces más y de continuo, empuja la
proliferación sin estímulo. Lo que la vuelve oncogénica son la dosis y la
temporización, no la secuencia.
\end{solution}

\begin{solution}{exo:b3:rna-regulation:9}
(a) Madre de laboratorio, padre silvestre: los huevos no llevan \omterm{def:b3:rna-regulation:ncrna}{ARNpi}
contra los elementos P (el linaje de la madre nunca se los encontró), los
elementos P paternos transponen libremente en la línea germinal de la
descendencia --- estéril (disgénica). (b) Madre silvestre: sus huevos
contienen \omterm{def:b3:rna-regulation:ncrna}{ARNpi} contra los elementos P, que los silencian en la
descendencia --- fértil. La asimetría es el depósito materno de \omterm{def:b3:rna-regulation:ncrna}{ARNpi},
una herencia de ARN y no de genes.
\end{solution}

\begin{solution}{exo:b3:rna-regulation:10}
La proteína integra su mensajero a lo largo de su propia vida
$\tau_{p}$. Los mensajeros viven $\tau_{m} = 1/(\gamma + \gamma' R)$;
durante $\tau_{p}$ la proteína ve por tanto unas $\tau_{p}/\tau_{m}$
vidas de mensajero independientes, cada una un suceso aleatorio de
nacimiento y muerte. La fluctuación relativa de una media sobre $N$
sucesos independientes cae como $1/\sqrt{N}$, de modo que el $\tau_{m}$
más corto bajo control de \omterm{def:b3:rna-regulation:ncrna}{microARN} --- a igual mensajero medio, lo que
exige un $k$ proporcionalmente mayor --- da una proteína más suave. La
célula paga el \omterm{def:b3:rna-regulation:ncrna}{microARN} y la transcripción adicional para comprar
precisión y velocidad: un nivel fijado por un equilibrio entre síntesis
rápida y degradación rápida es a la vez menos ruidoso y más rápido de
reajustar que el mismo nivel fijado por síntesis y degradación lentas.
\end{solution}

\begin{solution}{exo:b3:rna-regulation:11}
Gusano: la ARN-polimerasa dependiente de ARN fabrica \omterm{def:b3:rna-regulation:ncrna}{ARNip} nuevos a
partir del propio mensajero diana, de modo que la señal se renueva
mientras la diana se transcriba, sobrevive a la dilución por las
divisiones y, con un canal de ARN de doble hebra entre células, se
propaga a otros tejidos y a la línea germinal durante unas pocas
generaciones. Mamífero: sin amplificación, los \omterm{def:b3:rna-regulation:ncrna}{ARNip} se consumen y se
diluyen en cada división, actúan solo en las células que los recibieron y
se desvanecen en días en las células que se dividen. Consecuencia: un
fármaco de \omterm{def:b3:rna-regulation:ncrna}{ARNip} ha de entregarse dentro de cada célula diana, funciona
mejor en células que no se dividen (los hepatocitos, de ahí los fármacos
hepáticos) y exige dosis repetidas.
\end{solution}

\begin{solution}{exo:b3:rna-regulation:12}
(a) ARN funcional: destruir el transcrito después de fabricado (\omterm{def:b3:rna-regulation:ncrna}{ARNip},
oligonucleótido antisentido) da un fenotipo, y expresar el ARN desde un
transgén situado en otro punto del genoma lo rescata. (b) Transcripción
funcional: insertar una señal de poliadenilación temprana que detenga la
transcripción (o suprimir el promotor) da el fenotipo, mientras que
destruir el ARN después de transcrito no, y un transgén ectópico no
rescata --- lo que importa es el acto de transcribir (o el elemento de
ADN), como en muchos ARN asociados a potenciadores. (c) Ruido: ninguna
manipulación cambia nada, el ARN está a menos de una copia por célula, y
ni su secuencia ni su expresión están conservadas.
\end{solution}

\begin{solution}{pb:b3:rna-regulation:1}
\textbf{1.} $\gamma = \ln 2/40 = \qty{0.0173}{min^{-1}}$; $m^{*} =
3/0.0173 \approx 170$ mensajeros.
\textbf{2.} $\delta_{p} = \ln 2/180 = \qty{0.00385}{min^{-1}}$; $P^{*} =
\beta m^{*}/\delta_{p} = 2\times 173/0.00385 \approx 9.0\times 10^{4}$
proteínas.
\textbf{3.} Degradación $4\gamma = \qty{0.069}{min^{-1}}$; $m^{*} = 3/0.069
\approx 43$; constante de tiempo $1/0.069 = \qty{14}{min}$.
\textbf{4.} $\beta = 1$: $P^{*} = 43/0.00385 \approx 1.1\times 10^{4}$;
represión de $8$ veces ($4$ por la degradación, $2$ por la traducción).
\textbf{5.} Semitiempo del mensajero, $\ln 2/(4\gamma) = \qty{10}{min}$.
La proteína decae después con su propia semivida, \qty{3}{h}: lo que
limita el interruptor es la degradación de la proteína.
\textbf{6.} Sí: el mensajero baja solo hasta $1/4$ (aún se detecta con
facilidad) mientras que la proteína baja hasta $1/8$ y sigue bajando a
medida que se recambia la proteína vieja; la observación temprana de que
la proteína desaparecía ``mientras el ARNm permanecía'' reflejaba sobre
todo la mitad traduccional de la represión.
\textbf{7.} $10^{6}/250 = 4000$ moléculas por huevo.
\textbf{8.} $4000/550 \approx 7$ por célula al eclosionar. $4000/2^{n} < 1$
para $n \ge 12$ divisiones desde el cigoto --- unas tres divisiones
después de eclosionar.
\textbf{9.} La polimerasa copia el mensajero diana en ARN nuevo de doble
hebra, que se corta en \omterm{def:b3:rna-regulation:ncrna}{ARNip} secundarios: el desencadenante se regenera a
partir de la propia diana, de modo que sobrevive a la dilución y se
propaga (un canal pasa ARN de doble hebra entre células y a la línea
germinal). Se desvanece porque la amplificación necesita a la vez la
diana y un desencadenante primario; una vez desaparecido el ARN inyectado
la reserva secundaria se diluye generación tras generación y la línea
germinal se reinicia.
\textbf{10.} $5000/2^{n} < 100$: $2^{n} > 50$, $n = 6$ días (quedan $78$
copias).
\textbf{11.} Las células que no se dividen no diluyen el \omterm{def:b3:rna-regulation:rnai}{RISC}; el límite
es la vida química del \omterm{def:b3:rna-regulation:ncrna}{ARNip} (degradación por nucleasas, frenada por la
modificación química de las hebras) y la fuga lenta del depósito
endosómico que alimenta el citoplasma.
\textbf{12.} $2\times 10^{7}/\num{16384} \approx \num{1200}$
coincidencias por azar. Lo decidió la genética: la pérdida de
\emph{lin-14} da el fenotipo contrario al de la pérdida de \emph{lin-4},
los alelos de ganancia de función de \emph{lin-14} eran deleciones de
justamente los sitios de la 3$'$ UTR, y los siete sitios estaban
conservados.
\textbf{13.} Nivel $\propto$ semivida: subida de $50/10 = 5$ veces.
\textbf{14.} Caída de $1/0.05 = 20$ veces en la síntesis de ferritina.
\textbf{15.} La razón importación/almacenamiento sube $5\times 20 = 100$
veces del estado rico en hierro al pobre.
\textbf{16.} El agregado solo se ensambla cuando hay hierro citosólico
disponible, de modo que la conformación de la proteína es una lectura
directa de la concentración de hierro, sin receptor, sin hormona y sin
transcripción --- como un ribointerruptor, un metabolito detectado por la
molécula que controla.
\textbf{17.} La ferritina se traduce con independencia del hierro:
ferritina sérica muy alta, hierro sérico y depósitos normales (sin
sobrecarga). En el cristalino, que no tiene recambio, la ferritina se
acumula durante años y cristaliza --- una catarata precoz.
\textbf{18.} La horquilla es solo un punto de aterrizaje; el efecto es lo
que la proteína unida obstruya físicamente --- los ribosomas que rastrean
en el extremo 5$'$, una nucleasa en el extremo 3$'$.
\textbf{19.} Dos químicas cubren dos condiciones: el agregado de hierro y
azufre de IRP1 lo destruyen también los oxidantes y la falta de oxígeno,
de modo que IRP1 responde al estado redox, mientras que la degradación de
IRP2 dependiente de hierro responde al hierro mismo en todo el intervalo
fisiológico de oxígeno; la redundancia hace robusto al regulón, y las dos
informan de cosas distintas.
\textbf{20.} $12\times 48\times 33\times 2 = \num{38016}$. $1 - (1 -
50/\num{38016})^{50} = 1 - 0.99868^{50} \approx 1 - e^{-0.066} =
0.064$: alrededor del \qty{6}{\%}.
\textbf{21.} Las isoformas idénticas se unen de forma homofílica y se
repelen: las ramas de una misma neurona se evitan entre sí y se
despliegan, mientras que las ramas de neuronas distintas, que casi no
comparten isoformas, pueden cruzarse --- una identidad privada por
célula.
\textbf{22.} $2^{n}$; $\num{1024}$ para $n = 10$. En la realidad son
menos: la inclusión la deciden factores propios de cada tipo celular y
está coordinada entre exones, muchas combinaciones desplazan el marco de
lectura y las destruye la degradación mediada por codón sin sentido, y
los exones no son independientes.
\textbf{23.} \emph{Sxl} controla además la \omterm{def:b3:chromatin-epigenetics:x-inactivation}{compensación de dosis}: en las
hembras SXL bloquea la traducción de \emph{msl-2}, de modo que los dos
cromosomas X no se hipertranscriben; sin SXL un embrión XX ensambla el
complejo MSL sobre los dos cromosomas X y duplica su producción --- una
sobredosis letal de todos los genes ligados al X, antes incluso de que el
sexo esté en juego.
\textbf{24.} El promotor temprano solo dispara mientras está presente la
señal X:A; la SXL que fabrica fuerza el corte productivo del transcrito
del promotor de mantenimiento, activo en ambos sexos, de modo que una vez
existe SXL se sigue fabricando a sí misma y la señal deja de hacer falta
--- la misma lógica autosostenida de lector que recluta escritor que en
las marcas de \omterm{def:b3:chromatin-epigenetics:nucleosome}{cromatina}, aquí con una proteína y un corte.
\textbf{25.} Proteína LIN-14 reprimida $8$ veces; semitiempo del
interruptor del mensajero, \qty{10}{min}; un ARN sin amplificar baja de
una molécula por célula tras $12$ divisiones.
\end{solution}
